\documentclass[11pt]{article}

\usepackage[utf8]{inputenc}
\usepackage[T1]{fontenc}
\usepackage{geometry}
\usepackage{amsmath}
\usepackage{amssymb}
\usepackage{booktabs}
\usepackage{hyperref}
\usepackage{xurl}
\usepackage{xcolor}
\usepackage{listings}
\usepackage{enumitem}
\usepackage{parskip}
\usepackage{graphicx}
\graphicspath{{../images/}}

\geometry{margin=1in}

\lstset{
  language=Java,
  basicstyle=\ttfamily\small,
  keywordstyle=\color{blue},
  commentstyle=\color{gray},
  stringstyle=\color{teal},
  showstringspaces=false,
  columns=fullflexible,
  frame=single,
  framerule=0.2pt,
  breaklines=true
}

\setlist[itemize]{topsep=0.3em,itemsep=0.2em}
\setlist[enumerate]{topsep=0.3em,itemsep=0.2em}

% Simple per-section source line.
\newcommand{\notesources}[1]{\par\noindent{\small\color{gray}\textit{Sources:} \sloppy #1\par}}

% Unit-wise source strings for this subject (used by slides/notes).
% Path is relative to the lecture `latex/` folder (the compilation CWD).
% OOPS with Java (CSEG1044) — unit-wise sources
%
% These are short, student-facing reference strings intended to be used with:
%   - \slidesources{...}  (in beamer slides)
%   - \notesources{...}   (in lecture notes)
%
% Style inspiration: other subjects in My-Subjects/ use semicolon-separated sources
% and include an "accessed" date for web documentation.

\newcommand{\OOPSUnitISources}{Schildt, \textit{Java: A Beginner's Guide} (10e), McGraw-Hill (Java fundamentals + OOP basics).; Oracle Java Tutorials: Getting Started + Language Basics + Classes and Objects (accessed \today).; Java SE API docs (e.g., \texttt{String}, \texttt{Scanner}) (accessed \today).}

\newcommand{\OOPSUnitIISources}{Schildt, \textit{Java: The Complete Reference} (12e), McGraw-Hill (classes, inheritance, packages, interfaces).; Bloch, \textit{Effective Java} (3e), Addison-Wesley, 2018 (design choices; interfaces vs abstract classes).; Oracle Java Tutorials: Classes and Objects + Interfaces + Packages (accessed \today).; Java SE API docs (e.g., \texttt{Comparable}, \texttt{Runnable}) (accessed \today).}

\newcommand{\OOPSUnitIIISources}{Schildt, \textit{Java: The Complete Reference} (12e) (nested/inner classes, exceptions, threads, I/O).; Oracle Java Tutorials: Nested Classes + Exceptions + Concurrency + Basic I/O (accessed \today).; Java SE API docs (e.g., \texttt{Thread}, \texttt{AutoCloseable}) (accessed \today).}

\newcommand{\OOPSUnitIVSources}{Schildt, \textit{Java: The Complete Reference} (12e) (generics, lambdas, Swing, JDBC).; Bloch, \textit{Effective Java} (3e) (generics + lambdas).; Oracle Java Tutorials: Generics + Lambda Expressions + Swing + JDBC Basics (accessed \today).; Java SE API docs (e.g., \texttt{java.util.function}, \texttt{javax.swing}, \texttt{java.sql}) (accessed \today).}

\newcommand{\OOPSUnitVSources}{Schildt, \textit{Java: The Complete Reference} (12e) (Collections Framework + wrapper classes).; Oracle Java docs: Collections Framework overview (accessed \today).; Java SE API docs (\texttt{java.util} collections; wrapper classes in \texttt{java.lang}) (accessed \today).}

\newcommand{\OOPSUnitVISources}{Git SCM: \textit{Pro Git} (2e) (version control workflow), accessed \today.; GitHub Docs (repositories, issues, pull requests), accessed \today.; Oracle Java Tutorials: Swing + JDBC (accessed \today).; JUnit 5 User Guide (accessed \today).; Apache Maven Project: Getting Started (accessed \today).; Gradle User Manual: Getting Started (accessed \today).; UML class diagrams: UML 2.x overview/spec (accessed \today).}

\newcommand{\OOPSMiscWhyInterfacesSources}{Schildt, \textit{Java: The Complete Reference} (12e) (interfaces).; Bloch, \textit{Effective Java} (3e) (prefer interfaces where appropriate).; Oracle Java Tutorials: Interfaces (accessed \today).; Java SE API docs (e.g., \texttt{Comparable}, \texttt{Runnable}) (accessed \today).}



\title{Object Oriented Programming (CSEG1044)\\Unit 01 -- Lecture 01 Notes\\OOP Paradigm + Principles + Modelling}
\begin{document}
\maketitle
\tableofcontents

\section{Lecture Snapshot}
\begin{itemize}
  \item \textbf{Unit focus:} Introduction to OOPs and Java Fundamentals
  \item \textbf{Main new ideas:} Why OOP? (vs procedural), The 4 OOP principles (with intuition), How to model a real-world entity
  \item \textbf{Course thread:} Use OOP Paradigm + Principles + Modelling to make the first text-based version of Campus Resource Manager easier to reason about before the full object model arrives.
\end{itemize}
\notesources{\OOPSUnitISources}

\section{Lecture Overview}
This lecture starts the course with two questions:
\begin{enumerate}
  \item \textbf{Why do we use OOP?} (what problems it solves compared to purely procedural code)
  \item \textbf{How do we think in OOP?} (principles + modelling a real-world problem into classes and objects)
\end{enumerate}
We will connect ideas to simple, beginner-friendly Java examples (no advanced syntax yet).
\notesources{\OOPSUnitISources}

\section{Prerequisite Bridge}
Before reading the new material in depth, do a fast recall of the older ideas listed below.
\begin{itemize}
  \item \textbf{Programming basics}: Variables, loops, and functions from C still matter; OOP changes how they are organized. Main reference: Programming in C prerequisite.
  \item \textbf{Problem decomposition}: OOP still depends on breaking problems into clear responsibilities before coding. Main reference: Programming in C prerequisite.
\end{itemize}
This lecture only revisits those ideas briefly so that the main explanation can stay focused on the new concept sequence.
\notesources{\OOPSUnitISources}

\section{Learning Outcomes}
After this lecture, you should be able to:
\begin{itemize}
  \item Explain why OOP became popular for large software systems.
  \item Define: abstraction, encapsulation, inheritance, polymorphism (with a small example).
  \item Convert a short problem statement into a small class design (classes + responsibilities).
\end{itemize}

\section{Suggested Reading Order}
\begin{enumerate}
  \item Read the \textbf{Prerequisite Bridge} first and confirm each recalled idea.
  \item Study the central explanation in this order: \textit{Why OOP? (vs procedural)} $\rightarrow$ \textit{The 4 OOP principles (with intuition)} $\rightarrow$ \textit{How to model a real-world entity}.
  \item Trace the worked example line by line before checking short answers.
  \item Attempt the practice section without looking at the solution immediately.
  \item Finish with the exit questions and further-study suggestions to confirm retention.
\end{enumerate}
\notesources{\OOPSUnitISources}

\section{Key Terms (Quick)}
\begin{itemize}
  \item \textbf{Class}: A blueprint (type) describing data (fields) and operations (methods).
  \item \textbf{Object}: A runtime instance of a class (has identity + state + behavior).
  \item \textbf{State}: The data inside an object (values of its fields at a time).
  \item \textbf{Behavior}: What an object can do (methods).
  \item \textbf{Invariant}: A rule that must always remain true for an object (e.g., \texttt{balance >= 0}).
  \item \textbf{Interface (idea)}: A public view/contract of what operations are allowed (we will formalize this later in Java).
\end{itemize}

\section{Detailed Notes}
\subsection{Why OOP? (vs procedural)}
\textbf{Procedural programming} organizes software around \emph{procedures/functions} that operate on data.
This is great for small programs, but as the program grows, a few problems become common:
\begin{itemize}
  \item \textbf{Shared state becomes risky:} if many functions can modify the same data, it becomes hard to guarantee correctness.
  \item \textbf{Tight coupling:} one change in data structure often forces changes in many functions.
  \item \textbf{Poor ``real-world'' mapping:} business concepts (student, course, account, booking) are not first-class; you must manually manage relationships.
\end{itemize}

\textbf{Object-Oriented Programming (OOP)} organizes software around \emph{objects}.
An object keeps its data together with the operations that are allowed on that data.
That makes it easier to:
\begin{itemize}
  \item \textbf{Divide responsibilities} (each class owns a small part of the system).
  \item \textbf{Protect invariants} (only methods can change internal data).
  \item \textbf{Reuse designs} (inheritance, composition, and interfaces).
  \item \textbf{Build large systems} by combining small, understandable pieces.
\end{itemize}

\paragraph{Rule of thumb.} In OOP, \emph{data is not ``free'' to modify}. Most fields should be private and updated through methods that validate.

\subsection{The 4 OOP principles (with intuition)}
\subsubsection{Abstraction}
\textbf{Abstraction} means focusing on \emph{what} something does, not \emph{how} it does it.
We hide irrelevant details so that the user of a class sees a simple model.

\textbf{Example.} You know a \texttt{BankAccount} supports \texttt{deposit()} and \texttt{withdraw()}. You do not care about internal storage details.

\textbf{Pitfall.} If you expose internal details (public fields), users can bypass rules and break the model.

\subsubsection{Encapsulation}
\textbf{Encapsulation} means bundling data with methods and controlling access to protect invariants.
In Java this usually means \texttt{private} fields + public methods.

\textbf{Example.} If \texttt{balance} is private, only \texttt{deposit()} / \texttt{withdraw()} can change it, so we can ensure \texttt{balance >= 0}.

\textbf{Pitfall.} ``Getters/setters for everything'' is not real encapsulation if setters allow invalid states (e.g., \texttt{setBalance(-100)}).

\subsubsection{Inheritance}
\textbf{Inheritance} models an \textbf{is-a} relationship: a child class is a specialized form of a parent class.
It enables reuse of common code \emph{when the relationship is truly is-a}.

\textbf{Example.} \texttt{SavingAccount} is-a \texttt{BankAccount} (it has all account behaviors and adds interest calculation).

\textbf{Pitfall.} Do not use inheritance only to reuse code. If it is not is-a, prefer composition (has-a).

\subsubsection{Polymorphism}
\textbf{Polymorphism} means ``many forms''.
In Java, the most important form is \textbf{runtime polymorphism}:
when a parent reference points to a child object, the overridden method that runs is chosen at runtime.

\textbf{Example (idea).} A variable of type \texttt{BankAccount} can hold a \texttt{SavingAccount}; calling \texttt{withdraw()} executes the child's override if it exists.

\textbf{Pitfall.} Confusing \textbf{overloading} (same method name, different parameters) with \textbf{overriding} (same signature in subclass).

\subsection{How to model a real-world entity}
When you get a problem statement, you are doing \textbf{modelling}: converting real-world concepts into a class design.
One practical step-by-step method:
\begin{enumerate}
  \item \textbf{Read scenarios}: list 3--5 user actions (``deposit money'', ``book a room'').
  \item \textbf{Extract nouns}: candidate classes/entities (Account, Customer, Booking).
  \item \textbf{Extract verbs}: candidate methods/operations (deposit, book, cancel).
  \item \textbf{Define responsibilities}: what each class owns/does (avoid ``God'' classes).
  \item \textbf{Write invariants}: rules that must never break (balance non-negative, booking has start < end).
  \item \textbf{Validate with examples}: run your design on sample inputs mentally.
\end{enumerate}

\paragraph{CRC cards (quick idea).} For each class, write:
\textbf{C}lass name, \textbf{R}esponsibilities, and \textbf{C}ollaborators.
This simple technique prevents one class from doing everything.

\section{Running Project Connection}
Use OOP Paradigm + Principles + Modelling to make the first text-based version of Campus Resource Manager easier to reason about before the full object model arrives.
\begin{itemize}
  \item Use Campus Resource Manager as the running case study, or map the same idea to your own project domain if you are using another example.
  \item Ask where today's idea should live: model, service, DAO, UI, or team workflow.
  \item Use the lecture example as a smaller version of the same design choice you will make in the capstone.
\end{itemize}
\notesources{\OOPSUnitISources}

\section{Common Mistakes and Confusions}
\begin{itemize}
  \item Treating a class as only a data bucket with no behavior.
  \item Confusing one object instance with the class blueprint.
  \item Using inheritance before the responsibilities are understood.
\end{itemize}
\notesources{\OOPSUnitISources}

\section{Worked Example: \texttt{BankAccount} (Encapsulation + Invariants)}
\subsection*{Goal}
Model a bank account so that the invariant \texttt{balance >= 0} is always true.

\subsection*{Code}
\begin{lstlisting}
public class BankAccount {
    private final String accountNo;
    private double balance; // invariant: balance >= 0

    public BankAccount(String accountNo, double openingBalance) {
        if (accountNo == null || accountNo.isBlank())
            throw new IllegalArgumentException("accountNo required");
        if (openingBalance < 0)
            throw new IllegalArgumentException("openingBalance cannot be negative");

        this.accountNo = accountNo;
        this.balance = openingBalance;
    }

    public double getBalance() {
        return balance;
    }

    public void deposit(double amount) {
        if (amount <= 0) throw new IllegalArgumentException("amount must be positive");
        balance += amount;
    }

    public void withdraw(double amount) {
        if (amount <= 0) throw new IllegalArgumentException("amount must be positive");
        if (amount > balance) throw new IllegalArgumentException("insufficient balance");
        balance -= amount;
    }
}
\end{lstlisting}

\subsection*{What to notice}
\begin{itemize}
  \item Encapsulation: \texttt{balance} is \texttt{private}.
  \item Invariant protection: constructor and methods validate inputs.
  \item Abstraction: users do not directly ``set balance''; they request operations.
\end{itemize}
\notesources{\OOPSUnitISources}

\section{Demo / Observation Task}
Use this block as a short reproducible demo or observation task.
\begin{itemize}
  \item Reproduce the lecture's main example around \textit{Why OOP? (vs procedural)} once without looking at the solution first.
  \item Change one input, rule, or design choice in Campus Resource Manager and predict the result before checking it.
  \item Record one success case, one edge case, and one failure case so the idea becomes observable instead of purely theoretical.
\end{itemize}
\notesources{\OOPSUnitISources}

\section{Practice Check (with brief solutions)}
\begin{enumerate}
  \item \textbf{Model a \texttt{Student} entity.} Identify 3 fields and 3 methods.\par
  \textbf{Sample:}
  \begin{itemize}
    \item Fields: \texttt{rollNo}, \texttt{name}, \texttt{cgpa}
    \item Methods: \texttt{updateCgpa()}, \texttt{enrollCourse()}, \texttt{printTranscript()}
  \end{itemize}

  \item \textbf{Find the invariant.} For each, write one invariant:\par
  \textbf{(a)} Rectangle \textbf{(b)} Inventory item \textbf{(c)} Booking.\par
  \textbf{Sample:} Rectangle: width/height > 0; Booking: startDate < endDate.

  \item \textbf{Is-a or has-a?} Decide if it should be inheritance or composition:\par
  \textbf{(a)} Car and Engine \textbf{(b)} Dog and Animal.\par
  \textbf{Answer:} Car has-a Engine (composition). Dog is-a Animal (inheritance).
\end{enumerate}

\section{Self-Study Practice (no solutions)}
\begin{enumerate}
  \item Convert a real-world object (Library Book, ATM Card, Classroom) into a class with fields + methods + 1 invariant.
  \item Write a 5--7 line paragraph: ``Why OOP helps in large software projects'' using one example.
  \item Modify \texttt{BankAccount} to add \texttt{transferTo(BankAccount other, double amount)}.
\end{enumerate}

\section{Further Study}
\begin{itemize}
  \item Revisit the prerequisite references if any recall bullet still feels weak.
  \item Rework the lecture example with one changed assumption, input, or design choice.
  \item Read the textbook and documentation references listed in the source line for a second explanation of the same topic.
  \item Apply the same idea once inside Campus Resource Manager so that the concept moves from theory to design practice.
\end{itemize}
\notesources{\OOPSUnitISources}

\section{Exit Questions (with sample answers)}
\textbf{Q1.} Why do we keep fields private in OOP?\par
\textbf{Sample answer:} Private fields prevent external code from putting an object into an invalid state. Changes happen only through methods, so we can validate inputs and preserve invariants.

\vspace{0.6em}
\textbf{Q2.} What is the difference between abstraction and encapsulation?\par
\textbf{Sample answer:} Abstraction is about simplifying the model (what is exposed); encapsulation is about protecting the internal state and controlling access (how it is protected).

\vspace{0.6em}
\textbf{Q3.} Give one example of ``is-a'' and one example of ``has-a''.\par
\textbf{Sample answer:} Dog is-a Animal (inheritance). Car has-a Engine (composition).

\end{document}
